%% ch4_implementation.tex — Trimmed Implementation Chapter with Diagrams

\chapter{IMPLEMENTATION}

This chapter documents the realisation of the testbed from physical infrastructure
through to orchestrator deployment and experimental campaign setup. The agentic
reasoning architecture and individual agent designs are detailed in
Chapters~3 and~5 respectively; this chapter focuses on how those components
were physically instantiated, integrated, and operated.

%% ─────────────────────────────────────────────────────────────────────────────
\section{Physical Testbed Infrastructure}

The testbed spans five Linux virtual machines connected over a private NAT network
with static IP assignments. Table~\ref{tab:infrastructure} summarises each node.

\begin{table}[H]
\centering
\caption{Testbed Infrastructure Inventory}
\label{tab:infrastructure}
\begin{tabular}{p{2.6cm}p{3.4cm}p{3.0cm}p{2.8cm}}
\toprule
\textbf{Node} & \textbf{Role} & \textbf{Resources} & \textbf{OS} \\
\midrule
\texttt{kubemaster} & K8s control plane; orchestrator host & 4\,vCPU, 8\,GB RAM & Ubuntu 22.04 \\
\texttt{kube}       & K8s worker; UPF pod host             & 4\,vCPU, 8\,GB RAM & Ubuntu 22.04 \\
\texttt{core VM}    & Open5GS 5G Standalone core           & 2\,vCPU, 4\,GB RAM & Ubuntu 22.04 \\
\texttt{UERANSIM}   & gNB and UE emulation                 & 2\,vCPU, 4\,GB RAM & Ubuntu 22.04 \\
\texttt{MEC VM}     & Edge application and traffic sink    & 2\,vCPU, 4\,GB RAM & Ubuntu 22.04 \\
\bottomrule
\end{tabular}
\end{table}

The Flannel CNI plugin provides intra-cluster pod networking.
UPF pods run in host-network mode to access the physical GTP-U interface directly.

%% ─────────────────────────────────────────────────────────────────────────────
\section{5G Core and Slice Configuration}

The 5G Standalone core runs Open5GS with three dedicated SMF instances, one per
slice, each bound to a distinct UPF and non-overlapping IP pool.
Table~\ref{tab:slice_config} shows the complete slice-to-function mapping.

\begin{table}[H]
\centering
\caption{Slice Configuration: S-NSSAI to SMF and UPF Mapping}
\label{tab:slice_config}
\begin{tabular}{lllll}
\toprule
\textbf{Slice} & \textbf{S-NSSAI} & \textbf{DNN} & \textbf{UPF Subnet} & \textbf{SMF Service} \\
\midrule
eMBB  & SST=1, SD=0x000001 & \texttt{embb}  & 10.45.0.0/16 & \texttt{smfd-embb}  \\
URLLC & SST=2, SD=0x000002 & \texttt{urllc} & 10.46.0.0/16 & \texttt{smfd-urllc} \\
mMTC  & SST=3, SD=0x000003 & \texttt{mmtc}  & 10.47.0.0/16 & \texttt{smfd-mmtc}  \\
\bottomrule
\end{tabular}
\end{table}

The AMF YAML registers PLMN MCC=999, MNC=70 and binds the NGAP interface on SCTP
port 38412. Subscriber records in the UDM MongoDB store per-IMSI S-NSSAI entries,
enabling the NSSF to resolve slice selection at registration time.

Nine UE instances are emulated via UERANSIM, each establishing three concurrent
PDU sessions (one per slice), resulting in three \texttt{uesimtun} tunnel
interfaces per UE. Traffic generators bind to the correct interface, ensuring
packets enter the intended GTP-U tunnel without cross-slice leakage.

%% ─────────────────────────────────────────────────────────────────────────────
\section{User-Plane Isolation via HTB Traffic Shaping}

Per-slice bandwidth isolation is enforced at the UPF egress interface
(\texttt{ens33}) using Linux Traffic Control with a Hierarchical Token Bucket
(HTB) scheduler. Figure~\ref{fig:htb_tree} illustrates the class hierarchy.

\begin{figure}[H]
\centering
\begin{tikzpicture}[
  node distance=1.0cm and 1.8cm,
  box/.style={draw, rounded corners=3pt, minimum width=3.0cm,
              minimum height=0.7cm, align=center, font=\small},
  root/.style={box, fill=gray!15, thick},
  embb/.style={box, fill=blue!12, draw=blue!70!black},
  urllc/.style={box, fill=red!12, draw=red!70!black},
  mmtc/.style={box, fill=green!12, draw=green!60!black},
  arr/.style={-Latex, thick}
]
\node[root] (root) {\texttt{1:0} root HTB \\ \small 1000\,Mbit/s};
\node[root, below=1.2cm of root] (parent) {\texttt{1:1} parent class \\ \small rate=1000 ceil=1000};

\node[embb,  below left=1.4cm and 2.4cm of parent] (embb)
    {\texttt{1:10} eMBB \\ \small rate=1000 ceil=1000};
\node[urllc, below=1.4cm of parent] (urllc)
    {\texttt{1:20} URLLC \\ \small rate=50 ceil=200 \\ \textbf{prio\,1}};
\node[mmtc,  below right=1.4cm and 2.4cm of parent] (mmtc)
    {\texttt{1:30} mMTC \\ \small rate=10 ceil=50};

\draw[arr] (root)   -- (parent);
\draw[arr] (parent) -- (embb);
\draw[arr] (parent) -- (urllc);
\draw[arr] (parent) -- (mmtc);

% Filter labels
\node[font=\scriptsize, text=cembb,  below=0.15cm of embb]  {dst 10.45.0.0/16};
\node[font=\scriptsize, text=curllc, below=0.15cm of urllc] {dst 10.46.0.0/16};
\node[font=\scriptsize, text=cmmtc,  below=0.15cm of mmtc]  {dst 10.47.0.0/16};
\end{tikzpicture}
\caption{HTB shaping hierarchy on the UPF egress interface. URLLC carries
\texttt{prio~1} (strict priority), ensuring URLLC packets skip the queue
whenever aggregate utilisation approaches the physical ceiling.}
\label{fig:htb_tree}
\end{figure}

Packets are classified into HTB classes via U32 filters that match the
destination IP subnet of each slice's UPF address pool. The orchestrator
modifies the eMBB rate at runtime using \texttt{tc class change} over SSH;
the command takes effect within one kernel scheduling cycle
($<$1\,ms enforcement latency).

%% ─────────────────────────────────────────────────────────────────────────────
\section{QoS Monitoring}

A dedicated monitoring thread collects per-slice telemetry every 3\,s,
independent of the LLM inference loop. Table~\ref{tab:metrics} summarises the
metrics collected and their measurement method.

\begin{table}[H]
\centering
\caption{Per-Slice Metrics Collected at Each Monitoring Cycle}
\label{tab:metrics}
\begin{tabular}{p{2.8cm}p{4.2cm}p{5.8cm}}
\toprule
\textbf{Metric} & \textbf{Source} & \textbf{Purpose} \\
\midrule
URLLC p99 RTT & ICMP echo via \texttt{uesimtun1} & Primary SLA indicator; fed to Planning Agent \\
RTT slope     & Linear regression over rolling window & Trend input for C3 Chain-of-Thought reasoning \\
eMBB load fraction & TX byte delta / current rate limit & Root-cause indicator for C4 WLA decision \\
eMBB throughput & \texttt{uesimtun0} byte counter delta & Reported in dataset; plotted in results \\
mMTC PDR      & MQTT TX vs RX message count & Alert if delivery ratio $<$ 99.5\,\% \\
\bottomrule
\end{tabular}
\end{table}

All metrics are exported in Prometheus text format on port 9200 and visualised
via Grafana dashboards, providing a real-time audit trail throughout the campaign.

%% ─────────────────────────────────────────────────────────────────────────────
\section{Orchestrator Control Loop}

Figure~\ref{fig:control_loop} shows the closed-loop execution flow shared by
both the rule-based and agentic orchestrators. The two systems differ only in
the \textit{Decide} stage.

\begin{figure}[H]
\centering
\begin{tikzpicture}[
  node distance=0.6cm and 0.5cm,
  stage/.style={draw, rounded corners=4pt, minimum width=2.5cm,
                minimum height=0.8cm, align=center, font=\small, thick},
  obs/.style={stage, fill=gray!12},
  think/.style={stage, fill=blue!8, draw=blue!60!black},
  validate/.style={stage, fill=yellow!20, draw=yellow!60!black},
  act/.style={stage, fill=green!12, draw=green!60!black},
  reflect/.style={stage, fill=cyan!12, draw=cyan!60!black},
  arr/.style={-Latex, thick},
  lab/.style={font=\scriptsize, midway, above=1pt}
]
\node[obs]      (obs)      {Observe \\ \scriptsize Monitoring Agent};
\node[think,    right=1.3cm of obs]   (think)   {Think \\ \scriptsize Planning Agent};
\node[validate, right=1.3cm of think] (val)     {Validate \\ \scriptsize Gate};
\node[act,      right=1.3cm of val]   (act)     {Act \\ \scriptsize Execution Agent};
\node[reflect,  below=1.1cm of think]  (reflect) {Reflect \\ \scriptsize Memory Agent};

\draw[arr] (obs)   -- node[lab]{state}      (think);
\draw[arr] (think) -- node[lab]{decision}   (val);
\draw[arr] (val)   -- node[lab]{action}     (act);
\draw[arr] (act)   -- ++(0,-1.45) -| node[near start, below, font=\scriptsize]{outcome} (reflect);
\draw[arr] (reflect) -- node[lab, left, rotate=90, yshift=4pt]{\scriptsize memory} (think);
\draw[arr] (obs)  -- ++(0,-1.1) node[below left, font=\scriptsize]{$\Delta t = 3\,s$} -- ++(0,-0.3);

% Rule-based annotation
\node[font=\scriptsize, draw=gray!60!black, rounded corners=2pt, fill=gray!8,
      below=0.2cm of val, xshift=1.0cm]
      (rb) {\textit{Rule-based}: fixed $\theta$};
\node[font=\scriptsize, draw=blue!60!black, rounded corners=2pt, fill=blue!6,
      below=0.6cm of val, xshift=1.0cm]
      (ag) {\textit{Agentic}: LLM + CoT + WLA};
\end{tikzpicture}
\caption{Closed-loop orchestrator control flow. Both orchestrator types share the
Observe–Validate–Act–Reflect structure. The agentic system additionally invokes
the Memory Agent and produces a Chain-of-Thought reasoning trace.}
\label{fig:control_loop}
\end{figure}

In the rule-based system, the \textit{Think} stage is a binary threshold
comparison: if URLLC p99 RTT $>$ 20\,ms, throttle eMBB to 50\,Mbit/s;
otherwise restore to 1000\,Mbit/s. There is no root-cause analysis, memory, or
wrong-lever check. In the agentic system, the same stage invokes the local
\texttt{llama3.2:3b} model via Ollama, passes the enriched state including
memory context, and enforces the Chain-of-Thought schema and WLA gate before
forwarding a decision to the Execution Agent.

%% ─────────────────────────────────────────────────────────────────────────────
\section{Dataset Collection}

At each 3-second monitoring cycle, both orchestrators write a complete record to a
CSV log file partitioned by orchestrator type and load level.
Figure~\ref{fig:dataset_schema} illustrates the column structure of the collected
dataset.

\begin{figure}[H]
\centering
\begin{tcolorbox}[
  colback=gray!6, colframe=gray!40,
  title={\small\texttt{dataset\_agentic\_medium.csv} — Column Schema (6 datasets total)},
  fontupper=\ttfamily\scriptsize,
  left=4pt, right=4pt, top=4pt, bottom=4pt
]
\begin{tabular}{@{}ll@{\hspace{1.5cm}}ll@{}}
\multicolumn{2}{l}{\textbf{Network State}} & \multicolumn{2}{l}{\textbf{Agentic Reasoning}} \\
\texttt{urllc\_rtt\_ms}      & URLLC p99 RTT       & \texttt{root\_cause\_assessment} & CoT text field \\
\texttt{urllc\_rtt\_max\_ms} & URLLC max RTT       & \texttt{reasoning\_category}     & Classification \\
\texttt{embb\_mbps}          & Observed throughput & \texttt{lever\_validity\_score}  & WLA score (C4) \\
\texttt{embb\_rate\_mbit}    & Active tc rate      & \texttt{wla\_activated}          & WLA flag (C4)  \\
\texttt{mmtc\_msgs}          & mMTC message count  & \texttt{memory\_retrieval\_count}& Memory (C5) \\
\texttt{sla\_violated}       & SLA breach flag     & \texttt{memory\_reinforced}      & Reinforced (C5)\\
\texttt{orchestrator\_state} & Throttle/Normal     & \texttt{memory\_overridden}      & Override (C5) \\[4pt]
\multicolumn{2}{l}{\textbf{Decision}} & \multicolumn{2}{l}{\textbf{System}} \\
\texttt{action\_taken}       & Applied action      & \texttt{decision\_latency\_ms}   & LLM round-trip \\
\texttt{confidence}          & LLM confidence      & \texttt{tokens\_used}            & Token count \\
\texttt{candidate\_action}   & Pre-WLA action      & \texttt{node\_cpu\_percent}      & Host CPU load \\
\end{tabular}
\end{tcolorbox}
\caption{Column schema of the six collected datasets
(rule-based $\times$ \{low, medium, high\} and agentic $\times$ \{low, medium, high\}).
Each row is one 3-second control cycle. The agentic-specific columns
(\texttt{root\_cause\_assessment}, \texttt{lever\_validity\_score},
\texttt{wla\_activated}, \texttt{memory\_*}) are populated only in agentic runs;
rule-based runs carry equivalent synthetic labels for paired comparison.}
\label{fig:dataset_schema}
\end{figure}

Six CSV files are produced in total: one per combination of orchestrator type
(rule-based, agentic) and traffic load level (low, medium, high).
Each file contains approximately 9{,}000 rows, corresponding to a 7.5-hour
equivalent observation window per condition at the 3-second interval.

%% ─────────────────────────────────────────────────────────────────────────────
\section{Experimental Campaign Design}

The evaluation follows a $2 \times 3 \times 3$ factorial design:
2 orchestrator types $\times$ 3 load levels $\times$ 3 repeated trials,
yielding 18 total experimental runs. Figure~\ref{fig:exp_matrix} depicts
the campaign matrix.

\begin{figure}[H]
\centering
\begin{tikzpicture}[
  cell/.style={draw, minimum width=2.6cm, minimum height=0.75cm,
               align=center, font=\small},
  hdr/.style={cell, fill=gray!20, font=\small\bfseries},
  rb/.style={cell, fill=gray!8},
  ag/.style={cell, fill=blue!8, draw=blue!50!black}
]
% Headers
\node[hdr] (h0) at (0,0) {};
\node[hdr] (h1) at (2.7,0) {Low Load};
\node[hdr] (h2) at (5.4,0) {Medium Load};
\node[hdr] (h3) at (8.1,0) {High Load};
% Rule-based row
\node[hdr] (r1) at (0,-0.85)   {Rule-Based};
\node[rb]  at (2.7,-0.85)  {3 trials};
\node[rb]  at (5.4,-0.85)  {3 trials};
\node[rb]  at (8.1,-0.85)  {3 trials};
% Agentic row
\node[hdr] (r2) at (0,-1.70)  {Agentic};
\node[ag]  at (2.7,-1.70) {3 trials};
\node[ag]  at (5.4,-1.70) {3 trials};
\node[ag]  at (8.1,-1.70) {3 trials};
% Total
\node[font=\small\bfseries, right=0.4cm of h3] {$= 18$ runs};
\end{tikzpicture}
\caption{Experimental campaign matrix: 2 orchestrators $\times$ 3 traffic
load levels $\times$ 3 repeated trials = 18 total runs.}
\label{fig:exp_matrix}
\end{figure}

\noindent
Traffic load levels are defined as follows:

\begin{itemize}[noitemsep, topsep=4pt]
  \item \textbf{Low}: 360p HLS video (eMBB), 1\,Hz MQTT (mMTC), 1\,Hz ICMP (URLLC).
  \item \textbf{Medium}: 720p video, 2× MQTT publishers at 2\,Hz, 2\,Hz ICMP.
        Creates moderate inter-slice contention.
  \item \textbf{High}: 1080p video, 4× MQTT publishers at 4\,Hz, 4\,Hz ICMP.
        Generates sustained congestion reliably triggering SLA violations.
\end{itemize}

Table~\ref{tab:exp_params} lists the fixed parameters shared across all runs.

\begin{table}[H]
\centering
\caption{Experimental Parameters}
\label{tab:exp_params}
\begin{tabular}{ll}
\toprule
\textbf{Parameter} & \textbf{Value} \\
\midrule
URLLC SLA Threshold              & 20\,ms \\
Monitoring Interval ($\Delta t$) & 3\,s \\
eMBB Rate Bounds                 & 50 -- 1000\,Mbit/s \\
LLM Model                        & \texttt{llama3.2:3b} (Ollama, local CPU) \\
LLM Temperature                  & 0.2 \\
Memory Window                    & 10 decision records \\
Trials per Condition             & 3 \\
Total Experimental Runs          & 18 \\
\bottomrule
\end{tabular}
\end{table}

%% ─────────────────────────────────────────────────────────────────────────────
\section{Testbed Automation}

An automated shell script (\texttt{mec\_restart.sh}) restores the environment
to a known good state before each trial in five sequential phases:

\begin{enumerate}[noitemsep, topsep=4pt]
  \item \textbf{UPF Purge}: Force-delete UPF pods; wait for Kubernetes to
        respawn all three into \texttt{Running} state.
  \item \textbf{Core Restart}: SSH restart of AMF and all three SMF services;
        verify SCTP port 38412 binding.
  \item \textbf{RAN/UE Init}: Kill stale UERANSIM processes, delete residual
        tunnel interfaces, launch gNB and 9 UE instances, await PDU session
        establishment on all three slices.
  \item \textbf{Traffic Start}: Launch HLS, HTTP telemetry, and MQTT generators
        on their respective \texttt{uesimtun} interfaces.
  \item \textbf{Orchestrator Hook}: Start the Python orchestrator daemon; verify
        Prometheus exposure on port 9200.
\end{enumerate}

This procedure reduces environment setup from $\approx$25 minutes of manual
intervention to under 3 minutes of unattended execution, guaranteeing an
identical initial state for every trial.
